\subsection*{Properties of Relations}
\label{subsec:properties-of-relations}

\begin{definition}[Reflexive Relation]
\label{def:rel-reflexive}
Let $R$ be a binary relation on $A$.
The relation $R$ is \textbf{reflexive} if every element of $A$ is related to
itself:
\[
\forall a \in A,\quad a \mathrel{R} a.
\]
\end{definition}

\begin{remark*}[Standard quantified statement]
\[
\operatorname{ReflexiveRelation}(R,A)
\Longleftrightarrow
R\subseteq A\times A
\land
\forall a\in A,\; a\mathrel{R}a.
\]
\end{remark*}

\begin{remark*}[Definition predicate reading]
$R$ is reflexive on $A$ exactly when every element of $A$ has a loop to itself.
\end{remark*}

\begin{remark*}[Negated quantified statement]
\[
\neg\operatorname{ReflexiveRelation}(R,A)
\Longleftrightarrow
\exists a\in A,\; \lnot(a\mathrel{R}a).
\]
\end{remark*}

\begin{remark*}[Interpretation]
Reflexivity says that no element is excluded from relating to itself.
\end{remark*}

\begin{dependencies}
\begin{itemize}
  \item \hyperref[def:binary-relation]{Binary Relation}
\end{itemize}
\end{dependencies}

\begin{definition}[Irreflexive Relation]
\label{def:rel-irreflexive}
Let $R$ be a binary relation on $A$.
The relation $R$ is \textbf{irreflexive} if no element of $A$ is related to
itself:
\[
\forall a \in A,\quad \lnot(a \mathrel{R} a).
\]
\end{definition}

\begin{remark*}[Standard quantified statement]
\[
\operatorname{IrreflexiveRelation}(R,A)
\Longleftrightarrow
R\subseteq A\times A
\land
\forall a\in A,\; \lnot(a\mathrel{R}a).
\]
\end{remark*}

\begin{remark*}[Definition predicate reading]
$R$ is irreflexive on $A$ exactly when no element of $A$ is related to itself.
\end{remark*}

\begin{remark*}[Negated quantified statement]
\[
\neg\operatorname{IrreflexiveRelation}(R,A)
\Longleftrightarrow
\exists a\in A,\; a\mathrel{R}a.
\]
\end{remark*}

\begin{remark*}[Interpretation]
Irreflexivity forbids self-relation at every point of the carrier set.
\end{remark*}

\begin{dependencies}
\begin{itemize}
  \item \hyperref[def:binary-relation]{Binary Relation}
\end{itemize}
\end{dependencies}

\begin{definition}[Symmetric Relation]
\label{def:rel-symmetric}
Let $R$ be a binary relation on $A$.
The relation $R$ is \textbf{symmetric} if
\[
\forall a,b \in A,\quad
a \mathrel{R} b \Rightarrow b \mathrel{R} a.
\]
\end{definition}

\begin{remark*}[Standard quantified statement]
\[
\operatorname{SymmetricRelation}(R,A)
\Longleftrightarrow
R\subseteq A\times A
\land
\forall a,b\in A,\; a\mathrel{R}b \Rightarrow b\mathrel{R}a.
\]
\end{remark*}

\begin{remark*}[Definition predicate reading]
$R$ is symmetric on $A$ exactly when every related ordered pair can be reversed.
\end{remark*}

\begin{remark*}[Negated quantified statement]
\[
\neg\operatorname{SymmetricRelation}(R,A)
\Longleftrightarrow
\exists a,b\in A,\; a\mathrel{R}b \land \lnot(b\mathrel{R}a).
\]
\end{remark*}

\begin{remark*}[Interpretation]
Symmetry says that the relation has no preferred direction.
\end{remark*}

\begin{dependencies}
\begin{itemize}
  \item \hyperref[def:binary-relation]{Binary Relation}
\end{itemize}
\end{dependencies}

\begin{definition}[Antisymmetric Relation]
\label{def:rel-antisymmetric}
Let $R$ be a binary relation on $A$.
The relation $R$ is \textbf{antisymmetric} if
\[
\forall a,b \in A,\quad
a \mathrel{R} b \land b \mathrel{R} a \Rightarrow a=b.
\]
\end{definition}

\begin{remark*}[Standard quantified statement]
\[
\operatorname{AntisymmetricRelation}(R,A)
\Longleftrightarrow
R\subseteq A\times A
\land
\forall a,b\in A,\;
a\mathrel{R}b \land b\mathrel{R}a \Rightarrow a=b.
\]
\end{remark*}

\begin{remark*}[Definition predicate reading]
$R$ is antisymmetric on $A$ exactly when two-way relatedness can occur only
between equal elements.
\end{remark*}

\begin{remark*}[Negated quantified statement]
\[
\neg\operatorname{AntisymmetricRelation}(R,A)
\Longleftrightarrow
\exists a,b\in A,\;
a\mathrel{R}b \land b\mathrel{R}a \land a\neq b.
\]
\end{remark*}

\begin{remark*}[Interpretation]
Antisymmetry permits one-way comparison and self-comparison, but forbids
nontrivial cycles of length two.
\end{remark*}

\begin{dependencies}
\begin{itemize}
  \item \hyperref[def:binary-relation]{Binary Relation}
\end{itemize}
\end{dependencies}

\begin{definition}[Asymmetric Relation]
\label{def:rel-asymmetric}
Let $R$ be a binary relation on $A$.
The relation $R$ is \textbf{asymmetric} if
\[
\forall a,b \in A,\quad
a \mathrel{R} b \Rightarrow \lnot(b \mathrel{R} a).
\]
\end{definition}

\begin{remark*}[Standard quantified statement]
\[
\operatorname{AsymmetricRelation}(R,A)
\Longleftrightarrow
R\subseteq A\times A
\land
\forall a,b\in A,\; a\mathrel{R}b \Rightarrow \lnot(b\mathrel{R}a).
\]
\end{remark*}

\begin{remark*}[Definition predicate reading]
$R$ is asymmetric on $A$ exactly when every related ordered pair forbids its
reverse.
\end{remark*}

\begin{remark*}[Negated quantified statement]
\[
\neg\operatorname{AsymmetricRelation}(R,A)
\Longleftrightarrow
\exists a,b\in A,\; a\mathrel{R}b \land b\mathrel{R}a.
\]
\end{remark*}

\begin{remark*}[Interpretation]
Asymmetry is stricter than antisymmetry and implies irreflexivity.
\end{remark*}

\begin{dependencies}
\begin{itemize}
  \item \hyperref[def:binary-relation]{Binary Relation}
\end{itemize}
\end{dependencies}

\begin{definition}[Transitive Relation]
\label{def:rel-transitive}
Let $R$ be a binary relation on $A$.
The relation $R$ is \textbf{transitive} if related pairs can be chained:
\[
\forall a,b,c \in A,\quad
a \mathrel{R} b \land b \mathrel{R} c \Rightarrow a \mathrel{R} c.
\]
\end{definition}

\begin{remark*}[Standard quantified statement]
\[
\operatorname{TransitiveRelation}(R,A)
\Longleftrightarrow
R\subseteq A\times A
\land
\forall a,b,c\in A,\;
a\mathrel{R}b \land b\mathrel{R}c \Rightarrow a\mathrel{R}c.
\]
\end{remark*}

\begin{remark*}[Definition predicate reading]
$R$ is transitive on $A$ exactly when every two-step $R$-chain closes to a
one-step relation.
\end{remark*}

\begin{remark*}[Negated quantified statement]
\[
\neg\operatorname{TransitiveRelation}(R,A)
\Longleftrightarrow
\exists a,b,c\in A,\;
a\mathrel{R}b \land b\mathrel{R}c \land \lnot(a\mathrel{R}c).
\]
\end{remark*}

\begin{remark*}[Interpretation]
Transitivity is the chaining law behind order, equivalence, and divisibility.
\end{remark*}

\begin{dependencies}
\begin{itemize}
  \item \hyperref[def:binary-relation]{Binary Relation}
\end{itemize}
\end{dependencies}

\begin{definition}[Total Relation]
\label{def:rel-total}
Let $R$ be a binary relation on $A$.
The relation $R$ is \textbf{total} if every pair of elements is related in at
least one direction:
\[
\forall a,b \in A,\quad
a \mathrel{R} b \lor b \mathrel{R} a.
\]
\end{definition}

\begin{remark*}[Standard quantified statement]
\[
\operatorname{TotalRelation}(R,A)
\Longleftrightarrow
R\subseteq A\times A
\land
\forall a,b\in A,\; a\mathrel{R}b \lor b\mathrel{R}a.
\]
\end{remark*}

\begin{remark*}[Definition predicate reading]
$R$ is total on $A$ exactly when every pair of elements is comparable by $R$.
\end{remark*}

\begin{remark*}[Negated quantified statement]
\[
\neg\operatorname{TotalRelation}(R,A)
\Longleftrightarrow
\exists a,b\in A,\; \lnot(a\mathrel{R}b) \land \lnot(b\mathrel{R}a).
\]
\end{remark*}

\begin{remark*}[Interpretation]
Totality rules out incomparable pairs.
\end{remark*}

\begin{dependencies}
\begin{itemize}
  \item \hyperref[def:binary-relation]{Binary Relation}
\end{itemize}
\end{dependencies}
